\subsection{Finite-state constraints on defect tails}\label{sec:finite-state-tail}

Let $G$ be a finite directed graph with edge labels in
$\{-k,\ldots,k\}$, with distinct labels on the outgoing edges at each
vertex. Mark the edges labelled $k$ or $-k$. Let $S_k(G)\subset S_k$ be
the points whose eventual defect sequence labels an infinite path in $G$.
A strongly connected component is called marked if it contains an internal
marked edge. Write $\lambda_C$ for the Perron root of the adjacency matrix
of a marked component $C$, counting edge multiplicities.

\begin{theorem}[Finite-state defect tails]\label{thm:finite-state-tail}
If $G$ has no marked component, then $S_k(G)=\varnothing$.
Otherwise put $\lambda_*=\max_C\lambda_C$, where $C$ ranges over
the marked components. If $\lambda_*=1$, then $S_k(G)$ is countable.
If $\lambda_*>1$, it satisfies the three measure conclusions of
Corollary~\ref{cor:liminf-gauges}, with $S_k(G)$ in place of the named
set and with $b=\lambda_*$. In the zero case the null conclusion refers
to $S_k(G)$.
\end{theorem}

\begin{proof}
First consider a strongly connected marked graph with Perron root
$\lambda>1$. Take a positive right Perron vector $w$. The sum of
terminal weights over paths of length $m$ from vertex $v$ equals
$\lambda^m w_v$. Since the coordinates of $w$ are bounded above and
away from zero, the path count is comparable to $\lambda^m$, uniformly
over the initial vertex. The fixed initial-pair covering argument
therefore gives at most $C\lambda^{n-N}$ intervals of diameter $r_n$
for each fixed starting position, pair and vertex. This gives the
zero and finite-cover conclusions when the normalized lower limit is
zero or finite, respectively.

To obtain positive critical measure, give an edge $v\to v'$ probability
$w_{v'}/(\lambda w_v)$. The path-cylinder mass telescopes to
$w_{v_m}/(\lambda^m w_{v_0})\le C\lambda^{-m}$. Fix one marked edge.
From every vertex there is a path crossing it in at most some fixed
number $L_0$ of steps, with probability at least $p>0$. Conditioning
successively in blocks of length $L_0$ bounds the probability of
avoiding the edge for $m$ blocks by $(1-p)^m$. Applying this after each
starting time shows that almost every path crosses a marked edge
infinitely often.

Feed the edge labels into the defect recurrence, with the two
initial integers zero and with starting index $N>8k$ sufficiently
large. The construction in the proof of Theorem~\ref{thm:gauge} gives
a point $x=\lim M_j/Q_j$ with
$|Q_jx-M_j|\le 4kQ_j/Q_{j+1}<1/8$ for $j\ge N$.
Thus the $M_j$ are its nearest integers and the prescribed labels
are its defects. Infinitely many marked labels give $\Lam(x)=k/4$
by the defect asymptotic. The pushforward of the path probability
therefore gives full mass to $S_k(G)$.

If a set of diameter $r$ satisfies $r_{n+1}\le r<r_n$, the difference
between two of its coded values of $M_j$ is less than
$1/4+Q_jr$. For $j<n$ with $j\ge N$ this is less than one, and
at $j=n$ there are at most three possible adjacent integers.
Consequently the set meets at most three label cylinders of length
$n-N$. From a prescribed start vertex each label word specifies at
most one path. Thus its mass is at most $3C\lambda^{-(n-N)}$, which is
bounded by a fixed multiple of $h_\lambda(r)$ for the interpolated gauge
$h_\lambda(r_n)=\lambda^{-n}$. Singleton mass is zero by the same
cylinder estimate. The construction is in a real interval of length
less than $1/2$ if $N$ is large, so projection to the circle is
isometric there. The mass-distribution bound gives
$\mathcal H^{h_\lambda}(S_k(G))>0$.

For Borel measurability, intersect with the Borel level $S_k$. For each
start and vertex, the condition that every finite prefix is admissible
is a countable intersection of Borel conditions on the rounding
functions. An infinite path exists by finite branching. Take the
countable union over starts and the finite union over vertices.
For $y\in L$, the defect asymptotic gives $d_j(y)=0$ eventually.
Because the nearest-integer errors of $x\in S_k$ and of $y$ tend to
zero, their nearest integers add eventually, so $d_j(x+y)=d_j(x)$
eventually. Hence $S_k(G)$ is invariant under the dense group $L$
and is Lebesgue-null. If its positive critical measure were finite,
its normalized restriction would be invariant under this dense group
and, by continuity against continuous test functions, under the
entire circle. It would then be Haar measure, a contradiction.
For a nonempty open set $U$, the family $U+\ell$, $\ell\in L\bmod 1$,
covers the circle by density and has a finite subcover by compactness.
Invariance under those translations would make the whole level have
finite measure if $S_k(G)\cap U$ had finite measure. This proves local
infinitude. The comparison of adjacent scales in
Corollary~\ref{cor:liminf-gauges} now gives the three conclusions for
every gauge. The middle case combines the countable covering by
fixed initial-pair classes with local infinitude on the full level.

For a general finite graph, each infinite path eventually remains
in one strongly connected component, since the component graph is
acyclic. If its point is in $S_k$, its defects have absolute value
$k$ infinitely often. Thus the terminal component is marked, and
\[
 S_k(G)=\bigcup_{C\ {\rm marked}} S_k(C).
\]
For a marked component with $\lambda_C=1$, each vertex has exactly
one internal outgoing edge. Indeed, two outgoing edges at one vertex
give two return loops with distinct first edges. Traversing them in
the two orders gives two distinct loops of equal length. Concatenating
those equal-length loops contradicts the bounded path count at
Perron root one. A marked component is therefore a directed cycle
in this case. Each fixed start index, vertex and integer pair
then specifies a unique defect tail and at most one real point:
the recurrence determines all later $M_j$, and $M_j/Q_j\to x$.
This gives countability.

Finally, when $1<\lambda_C<\lambda_*$ and
$\rho_{\lambda_*}(g)<\infty$, a bounded subsequence at base $\lambda_*$
has normalized cost tending to zero at base $\lambda_C$ by the
factor $(\lambda_C/\lambda_*)^n$. Such components are null.
Countable components are also null. At least one component attains
$\lambda_*$, giving local infinitude or local non-$\sigma$-finiteness
in the corresponding regimes. The finite union of the component
covers gives the asserted global finite-cover conclusion.
\end{proof}

\begin{remark}[Two graph examples]
Forbidding successive marked labels gives the adjacency matrix
\[
 \begin{pmatrix}2k-1&2\\2k-1&0\end{pmatrix}.
\]
For $k\ge2$ its characteristic polynomial has opposite signs at $2$
and $2k+1$, so its Perron root is strictly between those values.
For a reducible example, take a disjoint unmarked vertex with the
$2k-1$ unmarked loops and a marked vertex with the two loops labelled
$\pm k$. For $k\ge2$ the larger adjacency root comes from the unmarked
vertex, while $S_k(G)$ has critical base $2$. Thus the root of the
whole adjacency matrix need not determine the gauge of an exact level.
\end{remark}

For comparison, the graph-directed similarity results of Mauldin
and Williams \cite[Theorems~3 and~4]{MW} use a spectral radius of weighted
similarity matrices and study the full construction. Here the scales
are supplied by $Q_n$, and the insistence on infinitely many marked
edges arises from exact-level membership. The local infinitude uses
the arithmetic translation group of the level.

\paragraph{Lean status.}
All three regimes of Theorem~\ref{thm:finite-state-tail} are now kernel-checked at Layer~A in
\texttt{EpiFormal/\allowbreak Finite\allowbreak State\allowbreak Complete.lean} ($12$ audited declarations
ending in \texttt{thm\_\allowbreak finite\_\allowbreak state\_\allowbreak tail\_\allowbreak complete} on $\R$
and \texttt{thm\_\allowbreak finite\_\allowbreak state\_\allowbreak tail\_\allowbreak circle\_\allowbreak complete}
on \texttt{UnitAddCircle}, building on \texttt{Finite\allowbreak State\allowbreak Gauge.lean},
\texttt{Finite\allowbreak State\allowbreak Tail.lean}, and
\texttt{Finite\allowbreak State\allowbreak Positive.lean}, with base axioms only and zero \texttt{sorryAx}).
Specifically, under the deterministic outgoing-label hypothesis of this subsection, the formalization proves
compactness and subset embedding of the deterministic path tree (\texttt{graphTree}),
direct finite-sum Frostman positivity (\texttt{mkMetric\_\allowbreak graphTree\_\allowbreak pos}),
unmarked-tail nullity via backward-reachability stabilization at $L=|V|$ (\texttt{mkMetric\_\allowbreak unmarkedTail\_\allowbreak null}),
$L$-translation invariance of exact graph tails, and the arbitrary-gauge $\liminf$ comparison;
the positive right Perron--Frobenius eigenvector $(w_v,\lambda_C)$ on each marked strongly connected component is taken as an explicit hypothesis rather than derived via Lean matrix-eigenvalue machinery.

